% Source: https://artofproblemsolving.com/community/c3h2787070_1997_manhattan_mathematical_olympiad_grade_6_problems_14
% Stable identity: aops-topic-2787070-q04
From the units column, $2N$ is 0 or 10. If $N=5$, its carry 1 would make $2E+1$ a multiple of 10 in the tens column, impossible. Hence $N=0$. The tens column then forces $E=5$ and carry 1, since $E\ne N$. Enumerate the remaining possible distinct digits in FORTY, compute the sum with TEN twice, and check the displayed digits. This exhaustive finite check returns exactly one assignment.
\begin{verbatim}
from itertools import permutations
answers=[]
for F,O,R,T,Y in permutations([1,2,3,4,6,7,8,9],5):
    forty=10000*F+1000*O+100*R+10*T+Y
    total=forty+2*(100*T+50)
    if not 10000<=total<=99999: continue
    S,I,X,t,y=map(int,str(total))
    if t!=T or y!=Y: continue
    if len({F,O,R,T,Y,5,0,S,I,X})==10:
        answers.append((forty,100*T+50,total))
assert answers==[(29786,850,31486)]
\end{verbatim}
Thus $\boxed{29786+850+850=31486}$.
